\documentclass[12pt]{article}
\usepackage[utf8]{inputenc}
\usepackage{amsmath}
\usepackage{amssymb}
\usepackage{graphicx}
\usepackage{hyperref}
\usepackage{geometry}
\geometry{a4paper, margin=1in}
\usepackage{xcolor}

\title{\textbf{Deconstructing RWKV \\ Part 1: Linear Attention and the WKV Operator}}
\author{Sovesh Mohapatra}
\date{\today}

\begin{document}

\maketitle

\section*{Introduction}
The Transformer architecture has dominated deep learning for nearly a decade. Its core mechanism—self-attention—enables rich contextual representations by comparing every token to every other token. But this comes at a steep cost: attention scales quadratically with sequence length, and autoregressive generation requires storing all past key-value pairs, leading to $O(L)$ memory that grows linearly with context.

Recurrent Neural Networks (RNNs) sit at the opposite extreme. They process sequences token-by-token, maintaining a fixed-size hidden state. This gives them $O(1)$ memory during inference and constant-time token generation. But RNNs suffer from vanishing gradients, cannot be parallelized during training, and typically underperform Transformers on complex tasks.

What if there existed an architecture that combined the best of both worlds? An architecture that trains in parallel like a Transformer but infers sequentially like an RNN?

Enter \textbf{RWKV (Receptance Weighted Key Value)}.

In this 3-part series, we will completely deconstruct RWKV from theory to pure PyTorch implementation, culminating in a benchmark comparing it against Transformers and LSTMs.

\section*{The Key Insight: Attention as Linear Recurrence}

The breakthrough behind RWKV comes from a deceptively simple observation: \textit{attention can be rewritten as a linear recurrence}.

Recall the standard self-attention formula:
\begin{equation}
    \text{Attention}(Q, K, V) = \text{softmax}\left(\frac{QK^T}{\sqrt{d}}\right)V
\end{equation}

For a single query position $t$, the output is:
\begin{equation}
    y_t = \sum_{i=1}^{t} \frac{\exp(q_t \cdot k_i / \sqrt{d})}{\sum_{j=1}^{t} \exp(q_t \cdot k_j / \sqrt{d})} v_i
\end{equation}

The denominator is a normalization term that depends on all previous keys. This is what makes attention quadratic—we must recompute it for every query.

Now consider a \textbf{linear attention} variant where we remove the softmax and use a different kernel:
\begin{equation}
    y_t = \frac{\sum_{i=1}^{t} \phi(q_t, k_i) v_i}{\sum_{i=1}^{t} \phi(q_t, k_i)}
\end{equation}

If we choose $\phi(q, k) = \exp(q) \cdot \exp(k)$ (an exponential kernel), something remarkable happens. The numerator becomes:
\begin{equation}
    \sum_{i=1}^{t} \exp(q_t) \cdot \exp(k_i) \cdot v_i = \exp(q_t) \cdot \sum_{i=1}^{t} \exp(k_i) \cdot v_i
\end{equation}

The term $\sum_{i=1}^{t} \exp(k_i) \cdot v_i$ is a \textbf{cumulative sum} that can be computed recursively:
\begin{equation}
    \text{kv}_t = \text{kv}_{t-1} + \exp(k_t) \cdot v_t
\end{equation}

This is the essence of RWKV: attention reformulated as a linear recurrence that can be computed in $O(1)$ memory.

\section*{The WKV Operator}

RWKV introduces the \textbf{WKV (Weighted Key-Value)} operator, which formalizes this recurrence with additional refinements:

\begin{equation}
    \text{wkv}_t = \frac{\sum_{i=1}^{t-1} \exp(-(t-1-i)w + k_i) v_i + \exp(u + k_t) v_t}{\sum_{i=1}^{t-1} \exp(-(t-1-i)w + k_i) + \exp(u + k_t)}
\end{equation}

Here:
\begin{itemize}
    \item $w$ is a learned \textbf{decay} parameter controlling how fast past information fades.
    \item $u$ is a learned \textbf{time-first} parameter giving extra weight to the current token.
    \item The exponential decay $\exp(-(t-1-i)w)$ implements a fading memory similar to an RNN's hidden state.
\end{itemize}

The numerator can be computed recursively:
\begin{equation}
    \text{num}_t = \exp(u + k_t) v_t + \exp(-w) \cdot \text{num}_{t-1}
\end{equation}

And similarly for the denominator:
\begin{equation}
    \text{den}_t = \exp(u + k_t) + \exp(-w) \cdot \text{den}_{t-1}
\end{equation}

The final output is simply:
\begin{equation}
    \text{wkv}_t = \frac{\text{num}_t}{\text{den}_t}
\end{equation}

\section*{Time Mixing vs. Channel Mixing}

RWKV decomposes the Transformer block into two distinct modules:

\subsection*{Time Mixing (TM)}
The Time Mixing module is RWKV's equivalent of self-attention. It computes the WKV recurrence to mix information across time. The key components are:
\begin{itemize}
    \item \textbf{Key projection}: $k_t = x_t W_k$
    \item \textbf{Value projection}: $v_t = x_t W_v$
    \item \textbf{Receptance}: $r_t = \sigma(x_t W_r)$ (a sigmoid gate, analogous to query)
    \item \textbf{WKV output}: $\text{wkv}_t$ computed via the recurrence above
    \item \textbf{Final output}: $o_t = r_t \odot \text{wkv}_t$ (element-wise multiplication)
\end{itemize}

The receptance gate is crucial—it controls how much of the WKV output flows through, similar to how a query determines which values to attend to.

\subsection*{Channel Mixing (CM)}
The Channel Mixing module is RWKV's equivalent of the feed-forward network (FFN). It uses a gated linear unit (GLU) variant with squared ReLU activation:
\begin{align}
    k_t &= x_t W_k \\
    r_t &= \sigma(x_t W_r) \\
    v_t &= \text{ReLU}(k_t)^2 W_v \\
    o_t &= r_t \odot v_t
\end{align}

The squared ReLU activation ($\text{ReLU}(x)^2$) is a key design choice that provides smoother gradients than standard ReLU while maintaining sparsity.

\section*{Why RWKV Matters}

RWKV offers several compelling advantages:

\begin{itemize}
    \item \textbf{Parallel Training}: The WKV recurrence can be unrolled and computed in parallel during training using a clever cumulative sum trick, achieving throughput comparable to Transformers.
    
    \item \textbf{O(1) Inference}: During autoregressive generation, RWKV maintains a constant-size hidden state (the accumulated numerator and denominator). No growing KV cache.
    
    \item \textbf{Linear Scaling}: Both training and inference scale linearly with sequence length, unlike Transformer's quadratic attention.
    
    \item \textbf{Transformer Performance}: RWKV achieves perplexity and accuracy comparable to Transformers on language modeling tasks, unlike standard RNNs.
\end{itemize}

\section*{Next Steps: From Math to Code}

The WKV operator may seem abstract, but implementing it in PyTorch is surprisingly straightforward. In \textbf{Part 2}, we will build the Time Mixing and Channel Mixing modules from scratch, implement both parallel training and recurrent inference modes, and assemble a complete RWKV model capable of autoregressive text generation.

\vspace{1em}
\noindent \textit{Stay tuned for the code drop as we reinvent RNNs for the Transformer era!}

\end{document}
